of Kummer theory the following fact, due to Grothendieck (SGA 1 XI, on the last page): if @@M0@@ is a proper scheme over a field @@M1@@, with @@M2@@, then there exists an isomorphism of @@M3@@-group schemes @@M4@@; consequently, if @@M5@@ is algebraically closed, we have @@M6@@.
\item If @@M7@@ has characteristic zero, then the kernel @@M8@@ is zero, since every finite unipotent group is then zero (SGA 3 XVII 1.7). One thus approximately recovers Kodaira's ``Vanishing theorem'' in the form given by Mumford [2].
\end{enumerate}

\sgaStableTarget{sga6.e13.thm.3-13}\textbf{Theorem 3.13.} \emph{Let @@M9@@ be a noetherian scheme and @@M10@@ a projective @@M11@@-scheme, endowed with a relatively ample invertible sheaf @@M12@@. We suppose that the geometric fibers of @@M13@@ are integral and of dimension @@M14@@. Then, for a family @@M15@@ of classes of invertible sheaves on the fibers of @@M16@@, the following conditions are equivalent}~:
\begin{enumerate}[label=(\roman*)]
\item \emph{@@M17@@ is a bounded family.}
\item \emph{In the Hilbert polynomials @@M18@@ of the @@M19@@, the coefficient @@M20@@ remains bounded and the coefficient @@M21@@ remains bounded below.}
\item \emph{As @@M22@@ ranges over @@M23@@, the degree @@M24@@ remains bounded and the intersection number @@M25@@ remains bounded below.}
\item \emph{There exists an integer @@M26@@ such that we have}
@@M27@@
\emph{in the sense of the order relation defined by ampleness\,; that is, for every @@M28@@, @@M29@@ and @@M30@@ are ample.}
\end{enumerate}

